\documentclass[../../../include/open-logic-section]{subfiles}

\begin{document}

\olfileid{sth}{card-arithmetic}{fix}
\olsection{$\aleph$નાં સ્થિરબિંદુઓ}

\olref[spine][]{chap}માં આપણે સૂચવ્યું હતું કે
પ્રતિસ્થાપનને વાજબી ઠરાવવાની જરૂર છે, કારણ કે તે
પદાનુક્રમને ઘણો ઊંચો બનાવે છે. થોડું ગણાંક-અંકગણિત
જોયા પછી આ ઊંચાઈનું એક ઉદાહરણ આપી શકીએ.

દેખીતી રીતે $0 < \aleph_0$, અને $1 < \aleph_1$,
અને $2 < \aleph_2$\ldots; ખરેખર દરેક પગલે કદનો
તફાવત વધુ \emph{મોટો} થતો જાય છે. તેથી દરેક
ક્રમવાચક સંખ્યા $\kappa$ માટે
$\kappa< \aleph_\kappa$ હશે એવું અનુમાન કરવાનું
મન થાય.

પરંતુ $\ZFC$માં આ અનુમાન \emph{ખોટું} છે.
ખરેખર \emph{$\aleph$નાં સ્થિરબિંદુઓ} અસ્તિત્વમાં
છે, એટલે એવા ગણાંકો $\kappa$ કે
$\kappa=\aleph_\kappa$.

\begin{prop}\ollabel{alephfixed}
$\aleph$નું સ્થિરબિંદુ છે.
\end{prop}

\begin{proof}
પુનરાવર્તનથી આ રીતે વ્યાખ્યાયિત કરો:
\begin{align*}
	\kappa_0 &= 0\\
	\kappa_{n+1} &= \aleph_{\kappa_n}\\
	\kappa&= \bigcup_{n < \omega}\kappa_n
\end{align*}
હવે \olref[cardinals][classing]{unioncardinalscardinal}
મુજબ $\kappa$ ગણાંક છે. વળી:
\[
	\kappa= \bigcup_{n < \omega} \kappa_{n+1} = 
	\bigcup_{n < \omega}\aleph_{\kappa_n} = 
	\bigcup_{\alpha < \kappa}\aleph_\alpha = \aleph_\kappa
\]
%	By construction, $\kappa$ is the least cardinal greater than each
%	$\kappa_n$. So $\aleph_\kappa$ is the least cardinal greater than
%	each $\aleph_{\kappa_n}$, and hence greater than each
%	$\kappa_{n+1}$. But equally $\kappa$ is the least cardinal greater
%	than each $\kappa_{n+1} = \aleph_{\kappa_n}$. So $\kappa=
%	\aleph_\kappa$.
\end{proof}

બૂલોઝે આપણે હમણાં રચેલા આ જ $\aleph$-સ્થિરબિંદુ
વિશે લેખ લખ્યો હતો. લેખની શરૂઆતમાં $\kappa$નું
અસ્તિત્વ નોંધ્યા પછી તેમણે કહ્યું:
\begin{quote}
	[ગણસિદ્ધાંતનો અગાઉનો પરિચય ન હોય એવા લોકોની
	દૃષ્ટિએ $\kappa$] \emph{ઘણી મોટી} સંખ્યા છે; એટલી
	મોટી કે મને લાગે છે, ઓછામાં ઓછી $\kappa$ વસ્તુઓ
	છે એવો દાવો કરતી કોઈ પણ સિદ્ધાંતની સત્યતા અંગે
	તે પ્રશ્ન ઊભો કરે છે.
	\citep[p.~257]{Boolos2000}
\end{quote}
લેખના અંતે તેમણે પૂછ્યું:
\begin{quote}
	[શું] આપણને એવી શંકા થાય છે કે વાર્તાની શરૂઆતમાં
	ગમે તેમ રહ્યું હોય, પરંતુ આટલે સુધી આવતાં ચક્રો
	ખાલી ફરી રહ્યાં છે અને હવે આપણે ખરેખર જે છે
	તેનું વર્ણન સાંભળી રહ્યાં નથી?
	\citep[p.~268]{Boolos2000}
\end{quote}
જો આપણે ખરેખર ``જે છે'' તેનાથી આગળ નીકળી
ગયા હોઈએ, તો દોષ સીધો પ્રતિસ્થાપન પર મૂકવો પડે.
કારણ કે આ સ્વયંસિદ્ધિ આપણો પુરાવો શક્ય બનાવે છે.
એવું હોય, તો કદાચ બૂલોઝે થોડાં દાયકાં પહેલાં
પ્રતિસ્થાપનનાં ``કોઈ અનિચ્છનીય'' પરિણામ નથી
એવો કરેલો દાવો ફરી વિચારવો પડે
(\olref[replacement][extrinsic]{sec} જુઓ).

પરંતુ $\kappa$નું અસ્તિત્વ ખરેખર એટલું ખરાબ છે?
અહીં રસેલનો \emph{ટ્રિસ્ટ્રમ શૅન્ડીનો વિરોધાભાસ}
વિચારીએ. ટ્રિસ્ટ્રમ શૅન્ડી પોતાના જીવનની દૈનિક
નોંધ લખે છે, પરંતુ એક દિવસની નોંધ લખતાં તેને એક
વર્ષ લાગે છે. દરેક પસાર થતા વર્ષે તે વધુ પાછળ
પડે છે: એક વર્ષ પછી તેણે ફક્ત એક દિવસ નોંધ્યો
હોય અને જીવેલા 364 દિવસ નોંધાયા ન હોય; બે વર્ષ
પછી ફક્ત બે દિવસ નોંધ્યા હોય અને 728 દિવસ બાકી
હોય; ત્રણ વર્ષ પછી ફક્ત ત્રણ દિવસ નોંધ્યા હોય
અને 1092 દિવસ બાકી હોય~\dots\footnote{અધિવર્ષો
ધ્યાનમાં લીધાં નથી.} છતાં ટ્રિસ્ટ્રમ
\emph{અમર} હોય, તો તે દરેક દિવસ નોંધશે, કારણ
કે પોતાના જીવનના $n$મા વર્ષે $n$મો દિવસ
નોંધશે. આમ ``સમયના અંતે'' તેની નોંધપોથી પૂર્ણ હશે.

હવે આ વિચારથી નીચેનો વિચાર એટલો જુદો શા માટે
છે: $\alpha$ એ $\aleph_\alpha$ કરતાં નાની છે,
અને સતત વધુ ને વધુ પાછળ પડે છે, પરંતુ
$\kappa$એ પહોંચતાં તેને સરખી પડે છે અને
$\kappa = \aleph_\kappa$ થાય છે?

આ મુદ્દો હાલ બાજુએ રાખીને અને $\ZFC$ સ્વીકારીએ
એમ ધારીને, સ્થિરબિંદુઓની કેટલીક વધુ રસપ્રદ
રચનાઓથી સમાપન કરીએ. આગળનાં ત્રણ પરિણામો સહજ
રીતે બતાવે છે કે કોઈ (બિનતુચ્છ) બિંદુએ પદાનુક્રમની
પહોળાઈ તેની ઊંચાઈ જેટલી થાય છે:

\begin{prop}\ollabel{bethfixed}
$\beth$નું સ્થિરબિંદુ છે, એટલે એવું $\kappa$ કે
$\kappa=\beth_\kappa$.
\end{prop}

\begin{proof}
\olref{alephfixed} જેવી જ રીત, ફક્ત ``$\aleph$''ને
બદલે ``$\beth$'' વાપરો.
\end{proof}

\begin{prop}\ollabel{stagesize}
$\card{V_{\omega+\alpha}} = \beth_{\alpha}$.
જો $\omega \ordtimes \omega \leq \alpha$ હોય, તો
$\card{V_\alpha} = \beth_\alpha$.
\end{prop}

\begin{proof}
પહેલો દાવો સરળ અનંતાતીત અનુમાનથી મળે છે. બીજો
દાવો એથી મળે છે કે $\omega \ordtimes \omega
\leq \alpha$ હોય, તો $\omega + \alpha = \alpha$.
આ સાબિત કરવા
\olref[ord-arithmetic][]{chap}ના ક્રમવાચક અંકગણિતનાં
તથ્યો વાપરીએ. પહેલાં નોંધો કે
$\omega \ordtimes \omega = \omega \ordtimes (1 \ordplus \omega) =
(\omega  \ordtimes 1) \ordplus (\omega\ordtimes\omega) = \omega
\ordplus (\omega \ordtimes \omega)$.
હવે $\omega \ordtimes \omega \leq \alpha$
હોય, એટલે કોઈ $\beta$ માટે
$\alpha = (\omega\ordtimes\omega) \ordplus \beta$,
તો $\omega \ordplus \alpha = \omega \ordplus
((\omega \ordtimes \omega) \ordplus \beta) =
(\omega \ordplus (\omega \ordtimes \omega)) \ordplus
\beta = (\omega \ordtimes \omega) \ordplus
\beta = \alpha$.
\end{proof}

\begin{cor}
એવું $\kappa$ છે કે $\card{V_\kappa} = \kappa$.
\end{cor}

\begin{proof}
\olref{bethfixed} મુજબ $\kappa$ને કોઈ
$\beth$-સ્થિરબિંદુ લો. દેખીતી રીતે
$\omega \ordtimes \omega < \kappa$.
તેથી \olref{stagesize}થી
$\card{V_\kappa} = \beth_\kappa= \kappa$.
\end{proof}

$V_\kappa$ની નીચે જેટલા તબક્કા છે એટલા જ
$V_\kappa$માં !!{element}s છે. તેથી સહજ રીતે
$V_\kappa$ જેટલું ઊંચું છે તેટલું પહોળું પણ છે.
આ ટ્રિસ્ટ્રમ-શૅન્ડીની વાત જેવું છે: એક તબક્કાથી
આગળના તબક્કે \emph{ઘાતગણ} લઈને જઈએ છીએ,
તેથી પદાનુક્રમ દરેક પગલે \emph{ઘણો} મોટો થાય છે.
પરંતુ ``અંતે'', એટલે $\kappa$ તબક્કે,
પદાનુક્રમની પહોળાઈ તેની ઊંચાઈ જેટલી થઈ જાય છે.

પ્રશ્ન થાય: \emph{પદાનુક્રમની પહોળાઈ તેની ઊંચાઈ
જેટલી કેટલી વાર થાય છે?} જવાબ છે:
\emph{જેટલી ક્રમવાચક સંખ્યાઓ છે એટલી વાર.}
પણ આ માટે થોડી સમજૂતી જોઈએ.

પદ $\tau$ને આમ વ્યાખ્યાયિત કરીએ.
કોઈ પણ $A$ માટે:
\begin{align*}
	\tau_0(A) & \defis \cardsucc{(\card{A})}\\
	\tau_{n+1}(A) & \defis \beth_{\tau_n(A)}\\
	\tau(A) & \defis \bigcup_{n < \omega}\tau_n(A)
\intertext{\olref{bethfixed}ની જેમ કોઈ પણ $A$ માટે
$\tau(A)$ એ $\beth$-સ્થિરબિંદુ છે અને
$\card{A} < \tau(A)$. હવે આ પુનરાવર્તી વ્યાખ્યા જુઓ:}
	W_0 &\defis 0\\
	W_{\alpha + 1} & \defis \tau(W_\alpha)\\
	W_\alpha & \defis \bigcup_{\beta < \alpha} W_\beta
	\text{, જ્યારે $\alpha$ સીમા ક્રમવાચક સંખ્યા હોય}
\end{align*}
\sourcecorrection{OLMD-143}{મૂળમાં
$\tau_0(A)=\card{A}$ છે. જો $\card{A}$ પહેલેથી
$\beth$નું સ્થિરબિંદુ હોય, તો આખું પુનરાવર્તન
ત્યાં જ અટકે અને મૂળનો ``$\card{A}<\tau(A)$''
દાવો ખોટો પડે. શરૂઆતમાં અનુગામી ગણાંક
$\cardsucc{(\card{A})}$ લેવાથી રચાયેલું
$\beth$-સ્થિરબિંદુ $\card{A}$ કરતાં કડક મોટું રહે છે.}
આ રચના બધી ક્રમવાચક સંખ્યાઓ માટે વ્યાખ્યાયિત
છે. તેથી સહજ રીતે $W$ એ ક્રમવાચક સંખ્યાઓમાંથી
$\beth$-સ્થિરબિંદુઓમાં ``!!a{injection}'' છે.
અને પહેલાંની જેમ જ કોઈ પણ $\alpha$ માટે
$V_{W_\alpha}$ જેટલું ઊંચું છે તેટલું પહોળું છે.

\end{document}
